% Irodov VI §6.3 original 148--158, printed p249, PDF p258.
% Note: source problem 155 prints "22 MeV", answer presumes 22 meV.
\atomictermsproblem{148}{249}
钠原子价电子处于主量子数$n=3$的状态，
且总角动量达到该主量子数所允许的最大值。
求原子磁矩模。
\exsolution
$n=3$允许最大$l=2$（D态），电子自旋$s=1/2$，
所以最大$J=l+s=5/2$，
谱项\({}^{2}D_{5/2}\)。
朗德因子
\[
 g_J=1+\frac{J(J+1)+s(s+1)-l(l+1)}
 {2J(J+1)}
 =\frac65.
\]
因此
\[
 \boxed{\mu=g_J\mu_B\sqrt{J(J+1)}
 =\frac65\mu_B\frac{\sqrt{35}}2
 =\frac{3\sqrt{35}}5\mu_B}.
\]
\excheck 题中$n=3$仍可有$3d$轨道，
最大$J$不是$3p$态的$3/2$。

\atomictermsproblem{149}{249}
某激发原子的电子构型为$1s^22s^22p\,3d$，
且处于该构型最大总角动量态。
求原子磁矩的大小。
\exsolution
闭合壳层角动量相互抵消；
未闭合$p,d$各一电子，
取最大$L=1+2=3$，
两个自旋平行$S=1$，最大$J=L+S=4$，
即\({}^{3}F_4\)。
朗德因子
\[
 g_J=1+\frac{20+2-12}{2(20)}=\frac54,
\]
所以
\[
 \boxed{\mu=\frac54\sqrt{20}\,\mu_B
 =\frac{5\sqrt5}2\,\mu_B}.
\]
\excheck $J$最大不代表$g=2$，
轨道与自旋磁旋比分别不同。

\atomictermsproblem{150}{249}
原子处于$S=3/2,L=2$态，
已知其总磁矩为零。
求总角动量模。
\exsolution
非零$J$时磁矩模
$\mu=g_J\mu_B\sqrt{J(J+1)}$。
允许$J=1/2,3/2,5/2,7/2$。
由$g_J=0$得
\[
 1+\frac{J(J+1)+15/4-6}{2J(J+1)}=0,
\quad3J(J+1)=\frac94.
\]
所以$J(J+1)=3/4$即$J=1/2$，
\[
 \boxed{M_J=\frac{\sqrt3}2\hbar}.
\]
\excheck 此处磁矩为零来自朗德因子$g_J=0$，
并非总角动量矢量模为零。

\atomictermsproblem{151}{249}
某原子自旋$S=2$，总角动量模为$\sqrt2\,\hbar$，
且磁矩为零。写出相应光谱项。
\exsolution
由$J(J+1)=2$得$J=1$。
磁矩零要求$g_J=0$，代入朗德公式：
\[
 1+\frac{2+S(S+1)-L(L+1)}{2(2)}=0.
\]
$S(S+1)=6$，整理得$L(L+1)=12$，
因此$L=3$（F项）。
多重度$2S+1=5$，
\[
 \boxed{{}^{5}F_1}.
\]
\excheck $S=2,L=3$允许$J=1$，
符合角动量三角不等式。

\atomictermsproblem{152}{249}
处于\({}^{2}P_{3/2}\)态的原子
置于磁感应强度$B=1.0\,\mathrm{kG}$的均匀磁场。
求总角动量矢量绕磁场的进动角频率。
\exsolution
对\({}^{2}P_{3/2}\)，$L=1,S=1/2,J=3/2$，
朗德因子
$g_J=1+[15/4+3/4-2]/[2(15/4)]=4/3$。
弱场进动频率
$\omega_L=g_J\mu_B B/\hbar$，
$1\,\mathrm{kG}=0.10\,\mathrm T$，
所以
\[
 \boxed{\omega_L=
 \frac{4\mu_B(0.10\,\mathrm T)}{3\hbar}
 \simeq1.17\times10^{10}\,\mathrm{rad/s}}.
\]
\excheck 高斯单位的$1$kG须换算为$0.1$T，
不可数值原样代入SI公式。

\atomictermsproblem{153}{249}
\({}^{2}P_{1/2}\)态原子在圆环电流轴上，
圆环半径$r=5.0\,\mathrm{cm}$，
电流$I=10\,\mathrm A$，原子离环心的轴向距离也为$r$。
求该环对原子的最大磁力。
\exsolution
\({}^{2}P_{1/2}\)有$L=1,S=1/2,J=1/2$，
$g_J=2/3$；
最大磁矩沿轴投影
$|\mu_z|=g_JJ\mu_B=\mu_B/3$。
SI制圆环轴向磁场
\[
 B(z)=\frac{\mu_0Ir^2}{2(r^2+z^2)^{3/2}},
 \qquad
 \left|\frac{dB}{dz}\right|_{z=r}
 =\frac{3\mu_0I}{2^{7/2}r^2}.
\]
因此
\[
 \boxed{F_{\max}=|\mu_z B'|
 =\frac{\mu_0I\mu_B}{2^{7/2}r^2}
 \simeq4.1\times10^{-27}\,\mathrm N}.
\]
\excheck 应使用磁矩的\emph{轴向投影}；
不能将磁矩模代入轴向梯度力。

\atomictermsproblem{154}{249}
基态氢原子在距长直导线$r=2.5\,\mathrm{cm}$处，
导线电流$I=10\,\mathrm A$。
求原子所受磁力的最大值。
\exsolution
氢基态为\({}^{2}S_{1/2}\)，$g_J=2,J=1/2$，
最大磁矩沿场方向投影
$|\mu_B|$。
长直导线场
$B(r)=\mu_0I/(2\pi r)$，
梯度模
$|dB/dr|=\mu_0I/(2\pi r^2)$。
因此
\[
 \boxed{F_{\max}=\mu_B\frac{\mu_0 I}{2\pi r^2}
 \simeq3.0\times10^{-26}\,\mathrm N}.
\]
\excheck 实际方向由原子的$m_J$取值决定，
本题给出的数值是力的最大模。

\atomictermsproblem{155}{249}
窄钒原子束处于基态\({}^{4}F_{3/2}\)，
经过长$l_1=5.0\,\mathrm{cm}$的横向非均匀磁场，
屏位于磁场后$l_2=15\,\mathrm{cm}$，
两个最外侧分束在屏上相距$\delta=2.0\,\mathrm{mm}$。
原印动能是$T=22\,\mathrm{MeV}$。
求磁场梯度，并说明与书末答案的差异。
\exsolution
该态$L=3,S=3/2,J=3/2$，
$g_J=1+[15/4+15/4-12]/[2(15/4)]=2/5$。
两最外侧$m_J=\pm J$磁矩投影差为
$\Delta\mu=2g_JJ\mu_B$。
设原子速度$v$，
磁区内时间$t_1=l_1/v$，
离开磁区至屏时间$t_2=l_2/v$，
则两束位移差
\[
 \delta=\frac{\Delta\mu}{M}B'
 \left(\frac{t_1^2}2+t_1t_2\right).
\]
用$T=Mv^2/2$消去质量和速率，
\[
 \boxed{B'=\frac{2T\delta}
 {g_JJ\mu_B l_1(l_1+2l_2)}}.
\]
若\emph{严格照原印}$T=22\,\mathrm{MeV}$，
\[
 \boxed{B'\simeq1.45\times10^{11}\,\mathrm{T/m}}.
\]
若能量单位实际应为$22\,\mathrm{meV}$，
则
\[
 \boxed{B'\simeq1.45\times10^2\,\mathrm{T/m}
 \simeq15\,\mathrm{kG/cm}},
\]
与原书书末给出的15kG/cm一致。
\excheck 原印MeV与书末数值隐含的meV相差
$10^9$倍；本译本明确列出两种结果，
不未经授权替换题面单位。

\atomictermsproblem{156}{249}
以下谱项在弱磁场中分裂为多少个子能级：
(a)\({}^{3}P_0\)；
(b)\({}^{2}F_{5/2}\)；
(c)\({}^{4}D_{1/2}\)？
\exsolution
线性塞曼能移
$\Delta E=g_J m_J\mu_BB$；
$m_J=-J,\ldots,J$。
(a)$J=0$，只有$m_J=0$，不劈裂；
(b)$J=5/2$且$g_J\ne0$，共有
$2J+1=6$个不同子能级；
(c)$L=2,S=3/2,J=1/2$，
朗德因子
\[
 g_J=1+\frac{3/4+15/4-6}{2(3/4)}=0,
\]
两磁态不发生一阶能量分离。
因此
\[
 \boxed{\text{(a)不分裂；(b)6个；(c)一阶不分裂}}.
\]
\excheck “(c)不分裂”不代表不存在两个$m_J$态，
而是它们在一阶弱场中能量重合。

\atomictermsproblem{157}{249}
弱磁场$B=2.50\,\mathrm{kG}$使谱项发生塞曼劈裂。
求(a)单重D项；(b)\({}^{3}F_4\)
的全劈裂能宽（eV）。
\exsolution
对总角动量$J$的子项，
端到端全劈裂
$\Delta E_{\rm tot}=2J|g_J|\mu_BB$。
(a)单重D项$L=J=2,S=0,g_J=1$，
\[
 \boxed{\Delta E_a=4\mu_B(0.25\,\mathrm T)
 \simeq5.79\times10^{-5}\,\mathrm{eV}}.
\]
(b)\({}^{3}F_4\)有$L=3,S=1,J=4$，
$g_J=5/4$，
\[
 \boxed{\Delta E_b=8(5/4)\mu_B(0.25\,\mathrm T)
 \simeq1.45\times10^{-4}\,\mathrm{eV}}.
\]
\excheck 这里全宽包含最大$m_J=+J$与最小$m_J=-J$
两端之差，不是相邻子级间隔。

\atomictermsproblem{158}{249}
在弱磁场中判断下列光谱跃迁呈正常还是反常塞曼效应：
(a)\({}^{1}P\to{}^{1}S\)；
(b)\({}^{2}D_{5/2}\to{}^{2}P_{3/2}\)；
(c)\({}^{3}D_1\to{}^{3}P_0\)；
(d)\({}^{5}I_5\to{}^{5}H_4\)。
\exsolution
线频移取决于
$\Delta\omega=(g_um_u-g_lm_l)\mu_BB/\hbar$，
仅当允许跃迁的频移组成等距三线群，
才表现为正常塞曼效应。
(a)单重态$S=0$、$g=1$，正常。
(b)两态分别$g_u=6/5$、
$g_l=4/3$，频移不是单一等间隔三线群，反常。
(c)下态$J_l=0$，上态$J_u=1$且
$g_u=1/2$，
只给三个等间隔频移，正常。
(d)\({}^{5}I_5\)有$L=6,S=2,J=5$，
\({}^{5}H_4\)有$L=5,S=2,J=4$，
分别计算朗德因子均为$9/10$，
所以允许$\Delta m=0,\pm1$只产生三线群，正常。
\[
 \boxed{\text{(a)正常；(b)反常；(c)正常；(d)正常}}.
\]
\excheck “正常”指观察到的等间隔三线型，
不要求所有项都恰有$g=1$。
